\hitsetup{
  statesecrets={公开},
  natclassifiedindex={TP242},
  intclassifiedindex={621.8},
  ctitlecover={面向航空制造的电连接器机器人自主抓取与装配系统研究},
  ctitle={面向航空制造的电连接器机器人自主抓取与装配系统研究},
  cxueke={工学},
  csubject={机械工程},
  caffil={机电工程学院},
  cauthor={待填写},
  csupervisor={待填写},
  cstudentid={待填写},
  cstudenttype={学术学位论文},
  etitle={Research on Robotic Autonomous Grasping and Assembly System for Electrical Connectors in Aviation Manufacturing},
  exueke={Engineering},
  esubject={Mechanical Engineering},
  eaffil={School of Mechatronics Engineering},
  eauthor={To Be Filled},
  esupervisor={To Be Filled},
  ckeywords={航空制造, 电连接器, 自主抓取, 柔顺装配, 机器人系统},
  ekeywords={aviation manufacturing, electrical connector, autonomous grasping, compliant assembly, robotic system}
}

\begin{cabstract}
本文面向航空制造场景下电连接器机器人自主抓取与装配任务，围绕复杂机舱环境中目标遮挡严重、柔性线缆干扰显著以及狭窄空间运动受限等问题，研究电连接器视觉感知、自主抓取规划与柔顺装配控制方法。论文以电连接器为核心对象，构建由目标识别与位姿估计、候选抓取生成与筛选、引导式避障路径规划以及装配控制组成的技术路线，为形成面向航空非结构化环境的自主装配系统提供方法基础。

在视觉感知方面，论文研究多实例连接器识别、标签身份区分、6D位姿估计及线缆掩膜提取方法；在自主抓取方面，研究面向装配任务的候选抓取生成、多约束抓取筛选以及避开柔性线缆障碍的运动规划策略；在装配执行方面，进一步为后续柔顺控制与系统集成提供状态输入与轨迹先验。现有内容已迁移至本LaTeX工程中，后续可在此基础上继续完善摘要、参考文献、图表与实验结果。
\end{cabstract}

\begin{eabstract}
This thesis focuses on robotic autonomous grasping and assembly of electrical connectors in aviation manufacturing scenarios. Aiming at severe occlusion, flexible cable interference, and motion constraints in narrow cabin environments, the work investigates visual perception, grasp planning, and compliant assembly methods for electrical connectors. The overall technical route includes target recognition and pose estimation, candidate grasp generation and selection, guided collision-free motion planning, and subsequent assembly-oriented control.

The current LaTeX project has been initialized according to the HIT thesis template, and the existing thesis content has been migrated into chapter files. It can be further refined by completing metadata, references, figures, tables, and experimental results.
\end{eabstract}
